\documentclass[11pt]{article}
\usepackage[a4paper,margin=22mm]{geometry}
\usepackage{amsmath,amssymb,amsthm}
\newtheorem{theorem}{Theorem}
\newtheorem{corollary}[theorem]{Corollary}
\title{GST all-dimensional spatial folds: a decisive geometric test of Hodge breaking versus algebraicity}
\date{23 October 2026}
\begin{document}
\maketitle
\begin{abstract}
We test the owner's two-outcome spatial hypothesis at every emerging geometric dimension: breaking the realized Hodge structure would provide a counterexample to the classical Hodge conjecture, whereas preserving it would establish the conjecture. The strongest noncircular test distinguishes full pure Hodge structure, the Hodge type of a fixed integral topological class, and canonical algebraic representability. A genuine family of smooth projective N-folds for every N at least two shows that a fixed rational class can lose its Hodge type under geometric folding even though the Hodge conjecture is true for every member in all codimensions. Actual topology-changing singular worlds can break purity, but lie outside the smooth projective statement. This falsifies the proposed decision rule, not the Hodge conjecture.
\end{abstract}
\section{Typed GST whole-world geometric experiment}
The core of GST General Space is dimensionless. Scheme dimension, complex structures, rational Hodge decompositions and topology only emerge at a geometric realization. In the existing GST graded algebra, signed charge and purity gap are separate. On each realization $Y$, distinguish its full purity condition $\mathsf P(Y)$, the typed rational Hodge condition $\mathsf H_p(Y,\beta)$ and actual canonical algebraicity $\mathsf A_p(Y,\beta)$:
\[
 \mathsf P(Y): H^m(Y,\mathbb Q)\text{ pure of weight }m\text{ for every }m,\quad
 \mathsf H_p(Y,\beta):\beta\in H^{2p}(Y,\mathbb Q)\cap H^{p,p}(Y),
\]
\[
 \mathsf A_p(Y,\beta):\beta=\operatorname{cl}^{can}_Y(Z)
 \text{ for some actual }Z\in CH^p(Y)_{\mathbb Q}.
\]
The rational Hodge conjecture is
$\forall Y\text{ smooth projective},p,\beta:\mathsf H_p(Y,\beta)\Rightarrow\mathsf A_p(Y,\beta)$.
We do not assume it for the experiment.

\section{Actual all-N smooth-projective Hodge alignment break with true algebraicity}
\begin{theorem}[The proposed breaking verdict has genuine false positives at every N at least two]
For every integer $N\ge2$, set $k=N-2$ and take actual smooth projective complex N-folds
\[
 Y_N^0=E_i\times E_i\times\mathbb P^k,\quad
 Y_N^1=E_i\times E_{2i}\times\mathbb P^k,\qquad
 E_\tau=\mathbb C/(\mathbb Z+\tau\mathbb Z).
 \tag{1}
\]
They are diffeomorphic under marked real torus coordinates, both have pure Hodge structures and the same entire Hodge polynomial
\[
 \boxed{\mathcal H_{Y_N^a}(u,v)=(1+u)^2(1+v)^2\sum_{j=0}^k(uv)^j.}\tag{2}
\]
The full rational Hodge conjecture is true on both in every codimension.
Nevertheless the very same marked integral cohomology class
\[
 \boxed{
 \beta_N=[(dx_1-dx_2)\wedge(dy_1-dy_2)]\smile h^k
 \in H^{2N-2}(Y_N^a,\mathbb Z),\quad
 h=c_1(\mathcal O_{\mathbb P^k}(1))
 } \tag{3}
\]
has type $(N-1,N-1)$ on $Y_N^0$ and is NOT of that type on $Y_N^1$.
In the aligned world $\beta_N$ is the genuine canonical algebraic class
of the elliptic diagonal times a point in $\mathbb P^k$,
up to global orientation sign.
\end{theorem}
\begin{proof}
Using marked real coordinates and holomorphic forms
$dz_1=dx_1+\tau\,dy_1$ and $dz_2=dx_2+\sigma\,dy_2$,
the exterior product is exactly
\[
 \boxed{
 (dx_1-dx_2)\wedge(dy_1-dy_2)\wedge dz_1\wedge dz_2
 =(\sigma-\tau)dx_1\wedge dy_1\wedge dx_2\wedge dy_2.
 } \tag{4}
\]
At $(\tau,\sigma)=(i,i)$ the scalar is zero; at $(i,2i)$ it equals $i$.
For a real two-class on the complex two-torus, vanishing of the
wedge with the nonzero holomorphic two-form is equivalent to its
$(1,1)$ Hodge type: the wedge detects $(0,2)$ and complex
conjugation detects $(2,0)$. Projective stabilization by $h^k$
is nonzero and transports these components to degree $2N-2$.
When the elliptic markings agree, the two-factor complex diagonal
is an actual algebraic divisor whose normal Thom class is (3)
at $k=0$; multiplying by the algebraic point class on
$\mathbb P^k$ yields the asserted cycle at every $k$.

Equation (2) is Kunneth applied to
$\mathcal H_E=(1+u)(1+v)$ and
$\mathcal H_{\mathbb P^k}=\sum_{j=0}^k(uv)^j$.
Both worlds are smooth projective and therefore their
whole Hodge structures are pure.

Both two-elliptic-factor bases are smooth projective
surfaces. Rational Hodge algebraicity on every such
surface follows in codimensions zero and two from
fundamental and point classes and in codimension one
from the Lefschetz $(1,1)$ theorem. The genuine
projective-bundle Chow and Hodge formulas for the
product with $\mathbb P^k$ commute with the canonical
Betti cycle map. Thus every rational Hodge class
in every codimension of either $Y_N^a$ is algebraic,
for every $N\ge2$, notwithstanding (4).
\end{proof}
\begin{corollary}[A proven false-positive verdict, not a counterexample to Hodge]
There exists a Hodge-type-breaking fold among genuine smooth
projective worlds in every dimension $N\ge2$ while the
rational Hodge conjecture remains true on BOTH worlds.
Therefore the implication
\[
 \boxed{\text{a fixed rational class loses Hodge type under a fold}
 \ \Longrightarrow\ \text{a counterexample to the Hodge conjecture}}
 \tag{5}
\]
is false as a proposed geometric decision rule.
The breaking class on $Y_N^1$ is not a Hodge class.
\end{corollary}

\section{Actual topological purity breaking in every dimension, but singular}
\begin{theorem}[Genuine nodal topology-fold mixed weight]
Let $C$ be the projective nodal curve made by identifying
two distinct points of an elliptic curve $E$.
Then for every $N\ge1$ the actual projective
$Z_N=C\times\mathbb P^{N-1}$ has dimension $N$
and
\[
 \boxed{\operatorname{Gr}^W_0H^1(Z_N,\mathbb Q)=\mathbb Q(0),
 \qquad\operatorname{Gr}^W_1H^1(Z_N,\mathbb Q)=H^1(E,\mathbb Q).}\tag{6}
\]
Its full cohomology is not pure in degree one,
but the variety is singular.
\end{theorem}
\begin{proof}
Normalization gives the short exact constant-sheaf
sequence
$0\to\mathbb Q_C\to\nu_*\mathbb Q_E\to\mathbb Q_{\rm node}\to0$.
The $H^0(E)\to H^0(\rm node)$ difference map is zero,
since the two branches carry the same constant.
The resulting mixed-Hodge exact sequence is
$0\to\mathbb Q(0)\to H^1(C,\mathbb Q)\to H^1(E,\mathbb Q)\to0$.
The projective factor does not affect $H^1$ and
does not remove the singular node.
\end{proof}

\section{Proof-audit verdict for the spatial binary decision law}
The first failure mode of the negative branch is precise:
a marked cohomology class can cease to satisfy $(p,p)$
when complex geometry is folded, \emph{even if the classical
Hodge conjecture holds on every world involved}.
The second failure mode is precise:
a genuine loss of \emph{whole Hodge purity} can be
created by a topology fold, but its output is singular
and therefore not an admissible smooth-projective
counterexample. Within the smooth projective category,
purity itself holds unconditionally by Hodge theory.

For the positive branch, preservation of full purity
on all smooth projective worlds is automatic; the
additional implication from this preservation to
canonical algebraic representability is exactly the
open conclusion that must be derived, not silently
identified with purity. These facts do not disprove
the Hodge conjecture or rule out new GST geometry.
They give a decisive all-dimensional mathematical
diagnosis of what a valid fold experiment has to
measure before its output can decide the conjecture.

\section{Verification}
An independent Wolfram exact determinant and full-Hodge
polynomial computation checks the aligned and misaligned
parameters and the projective stabilization for $N=2$
through $26$. The universal all-$N$ assertion is proved
by the product, Hodge and projective-bundle formulas.
The existing Lean charge/purity-gap laws already track
the distinct logical meanings of the two observed breaks.
\end{document}